\section{The electric rest-frame inequality and finitely many ends}
\label{chext:sec:charged-rest-extensions}

The one-ended charged theorem also gives the mass lower bound on the
closed area--charge branch.  We first record its purely electric
rest-frame consequence.  We then treat additional asymptotically flat
ends by removing their distant tails and regarding the resulting
spheres as further inner boundary components.  This second argument
uses only the electric rest-frame consequence.

\begin{corollary}[Electric rest-frame inequality]
\label{chext:cor:charged-electric-rest}
Let $(\Omega^3,g,K,\mathcal E)$ be smooth and orientable, with $\Omega$
connected and complete with its nonempty compact smooth boundary $S$
included, and with exactly one asymptotically flat end.  The boundary
may have finitely many components.  Define
\[
 16\pi\mu=R_g+(\operatorname{tr}_gK)^2-|K|_g^2,
 \qquad
 8\pi J=\operatorname{div}_g
       \bigl(K-(\operatorname{tr}_gK)g\bigr).
\]
Assume that $\mu,|J|_g$ are integrable, that the ADM limits exist, and
that, for some $q>1/2$, the end satisfies
\[
 g-\delta=O_2(r^{-q}),\qquad
 K=O_1(r^{-1-q}),\qquad \mathcal E=O_1(r^{-2}).
\]
There is no magnetic field.  Suppose throughout $\Omega$ that
\begin{equation}\label{chext:eq:charged-electric-matter-dec}
 \operatorname{div}_g\mathcal E=0,
 \qquad
 \mu-\frac{|\mathcal E|_g^2}{8\pi}\ge |J|_g,
\end{equation}
and on $S$, with normal into $\Omega$, that
$H+\operatorname{tr}_S K\le0$.
Let $a=a_g(S)$ be the minimum enclosing area over all smooth filled
cuts of the entire boundary, and put
\[
 r_a=\sqrt{\frac{a}{4\pi}},\qquad
 Q=\frac1{4\pi}\int_S
           \langle\mathcal E,\nu\rangle_g\,dA_g.
\]
If the ADM momentum vanishes and $r_a\ge |Q|$, then the ADM energy
$m$ is positive and
\begin{equation}\label{chext:eq:charged-electric-rest-bound}
 m\ge\frac12\left(r_a+\frac{Q^2}{r_a}\right).
\end{equation}
\end{corollary}

\begin{proof}
The electromagnetic momentum is zero when the magnetic field is
zero.  Thus \eqref{chext:eq:charged-electric-matter-dec} is precisely the
non-electromagnetic matter dominant energy condition in
Theorem~\ref{ch:intro:main}.  Closedness of
$\iota_{\mathcal E}dV_g$ identifies the displayed boundary charge
with the charge at infinity.  All other one-ended hypotheses are
unchanged.  That theorem gives $m>0$, $m\ge|Q|$, and
\[
 r_a\le m+\sqrt{m^2-Q^2}.
\]
The enclosing area is positive.  Moreover,
$m-\sqrt{m^2-Q^2}\le|Q|\le r_a$.  Consequently
\[
 \bigl(r_a-m+\sqrt{m^2-Q^2}\bigr)
 \bigl(r_a-m-\sqrt{m^2-Q^2}\bigr)\le0.
\]
Expanding and dividing by $2r_a>0$ proves
\eqref{chext:eq:charged-electric-rest-bound}.  This calculation includes
$r_a=|Q|$; no limiting argument or equality classification is needed.
\end{proof}

\begin{corollary}[Outer area-minimizing boundary]
\label{chext:cor:charged-electric-boundary-area}
Under the hypotheses of Corollary~\ref{chext:cor:charged-electric-rest},
suppose that every filled enclosing cut of the whole boundary has area
at least $A=\operatorname{Area}_g(S)$.  If $A\ge4\pi Q^2$, then
\[
 m\ge\sqrt{\frac{A}{16\pi}}+Q^2\sqrt{\frac{\pi}{A}}.
\]
In particular this applies to maximal purely electric source-free
electrovacuum data with such a boundary and zero ADM momentum.
\end{corollary}

\begin{proof}
The boundary itself is an admissible cut, so the area assumption gives
$a_g(S)=A$.  Apply Corollary~\ref{chext:cor:charged-electric-rest}.
For the stated electrovacuum specialization,
$\operatorname{tr}_gK=0$, $\operatorname{div}_gK=0$, and
$R_g=|K|_g^2+2|\mathcal E|_g^2$ imply
$\mu=|\mathcal E|_g^2/(8\pi)$ and $J=0$, hence
\eqref{chext:eq:charged-electric-matter-dec}.
\end{proof}

\subsection{The area associated with a designated end}

Let $\Omega$ instead have finitely many asymptotically flat ends
$\mathscr E_0,\ldots,\mathscr E_k$, with $\mathscr E_0$ designated,
and compact smooth boundary $S$.  We allow $S=\varnothing$ only when
$k\ge1$.  A designated-end cut is the full compact smooth manifold
boundary $\Gamma=\partial Y$ of a smooth closed codimension-zero
exterior $Y\subset\Omega$.  Its intrinsic interior
$W=\operatorname{int}Y$ must be connected, lie in
$\operatorname{int}\Omega$, contain a distant tail of
$\mathscr E_0$, and exclude distant tails of every other end and
the entire $S$.  Here $\partial Y$ means the boundary of $Y$ as a
manifold, not its relative topological boundary in $\Omega$;
in particular, every portion coinciding with $S$ is counted.
Bounded complementary pockets are assigned to the excluded side.
The cut need not be connected, and its individual components need
not separately separate $\Omega$.

Define
\begin{equation}\label{chext:eq:charged-designated-area}
 a_0=\inf_{\Gamma}\operatorname{Area}_g(\Gamma),\qquad
 r_0=\sqrt{\frac{a_0}{4\pi}},
\end{equation}
where the infimum is over these full designated-end cuts.  This is a
single enclosing-area infimum, not the sum of separate end or boundary
component infima.  A distant sphere in $\mathscr E_0$ is a competitor,
so $a_0<\infty$.

\begin{theorem}[Electric rest-frame extension to finitely many ends]
\label{chext:thm:charged-electric-finite-ends}
Let $(\Omega^3,g,K,\mathcal E)$ be smooth, connected and orientable,
complete with its compact smooth boundary $S$ included, and have the
finitely many asymptotically flat ends just described.  Assume
\eqref{chext:eq:charged-electric-matter-dec} everywhere, integrable
$\mu,|J|_g$, and $H+\operatorname{tr}_S K\le0$ with the normal into
$\Omega$ on every component of $S$.  Assume the end falloff of
Corollary~\ref{chext:cor:charged-electric-rest} on each end, with a common
$q>1/2$, and finite ADM limits at the designated end.  There is no
magnetic field.  Let
\[
 Q_0=\frac1{4\pi}\int_{S_R^0}
             \langle\mathcal E,\nu_R\rangle_g\,dA_g
\]
be its conserved outward charge.  Suppose its ADM momentum is zero
and that $0<a_0$ and $r_0\ge|Q_0|$, with $a_0$ as in
\eqref{chext:eq:charged-designated-area}.  Then its ADM energy $m_0$ is
positive and
\begin{equation}\label{chext:eq:charged-finite-ends-bound}
 m_0\ge\frac12\left(r_0+\frac{Q_0^2}{r_0}\right).
\end{equation}
No momentum condition is imposed at the other ends.
\end{theorem}

\begin{proof}
For $k=0$, this is Corollary~\ref{chext:cor:charged-electric-rest}.  Suppose
$k\ge1$.  In each unwanted end $\mathscr E_j$, choose a coordinate
sphere $S_{R_j}^j$ sufficiently far out, and remove its open outside
tail.  Denote the retained manifold by $\Omega_{\mathbf R}$.
The radii can be chosen beyond a connected compact core: fix paths
in $\Omega$ connecting the finitely many inner end collars and the
components of $S$, and put these paths inside that core before
increasing the radii.  Each retained end annulus is connected and
attached to the core.  Thus $\Omega_{\mathbf R}$ is connected, has
exactly the designated end, and has nonempty compact smooth boundary
\[
 S_{\mathbf R}=S\ \sqcup\ \bigsqcup_{j=1}^kS_{R_j}^j.
\]
It is a closed subset of the original complete manifold.  A Cauchy
sequence for its intrinsic length metric therefore has a limit in
that subset; smooth boundary half-ball charts give convergence in
the intrinsic metric as well.  This proves completeness with its
boundary included.

On an added sphere, the normal pointing into the retained exterior
is the negative of the normal pointing toward that unwanted infinity.
The asymptotic estimates give
\[
 H=-\frac2{R_j}+O(R_j^{-1-q}),\qquad
 \operatorname{tr}_{S_{R_j}^j}K=O(R_j^{-1-q}).
\]
Hence its future outward expansion, in the retained-exterior
orientation, is strictly negative for large $R_j$.  The original
boundary retains its weakly trapped sign.  All fields are merely
restricted, so the charged DEC and divergence equation are unchanged.
No extension of the fields over a cap is involved.  The designated
end and its ADM energy and momentum are also unchanged.

Put $\mathcal F=\iota_{\mathcal E}dV_g$.  The equation
$d\mathcal F=0$ and Stokes' theorem on the region between
$S_{\mathbf R}$ and a distant sphere of the designated end give
\begin{equation}\label{chext:eq:charged-tail-flux}
 \frac1{4\pi}\int_{S_{\mathbf R}}\mathcal F=Q_0,
\end{equation}
where every inner component is oriented into the retained exterior.
In particular, if $Q_j$ is the outward charge of the $j$th unwanted
end and $Q(S)$ is the charge of the original boundary with normal
into $\Omega$, then $Q_0=Q(S)-\sum_{j=1}^kQ_j$.
The charge in \eqref{chext:eq:charged-finite-ends-bound} is therefore the
designated end charge, not generally $Q(S)$.

Let $a_{\mathbf R}=a_g(S_{\mathbf R})$ be the minimum enclosing
area in the retained one-ended manifold.  Every smooth filled cut
of $S_{\mathbf R}$ also gives a designated-end cut in $\Omega$:
keep its chosen-end exterior, and assign phase zero to all removed
tails.  This introduces no new frontier because the original cut
already encloses every added inner sphere.  If the cut coincides
with any added sphere, that sphere is the same smooth surface in
the original end and its full area remains counted.  The same
observation applies to coincident portions of $S$.  Selecting the
component toward the designated end and filling bounded pockets
keeps the whole-boundary convention and can only remove boundary
components.  Consequently
\begin{equation}\label{chext:eq:charged-area-under-tail-removal}
 a_{\mathbf R}\ge a_0,
 \qquad
 r_{\mathbf R}:=\sqrt{\frac{a_{\mathbf R}}{4\pi}}
       \ge r_0\ge|Q_0|.
\end{equation}
This argument compares the smooth full-cut classes directly; it
does not require attainment of either infimum or a limiting
statement as $\mathbf R\to\infty$.

All hypotheses of Corollary~\ref{chext:cor:charged-electric-rest} now hold
on $\Omega_{\mathbf R}$, with charge $Q_0$ by
\eqref{chext:eq:charged-tail-flux}.  Hence
\[
 m_0\ge\frac12\left(r_{\mathbf R}
                        +\frac{Q_0^2}{r_{\mathbf R}}\right).
\]
The function $r\mapsto(r+Q_0^2/r)/2$ is nondecreasing for
$r\ge|Q_0|$ (and $r>0$).  Combining this fact with
\eqref{chext:eq:charged-area-under-tail-removal} proves
\eqref{chext:eq:charged-finite-ends-bound}, including the endpoint
$r_0=|Q_0|$.
\end{proof}

The finite-end theorem is explicitly a purely electric result in the
rest frame of the designated end.  The use of total full-cut area
is essential: taking separate component areas, or applying a sum
of end-energy estimates, would not give the comparison
\eqref{chext:eq:charged-area-under-tail-removal}.
